% C.3au. Full author and independent root proof/source audits passed.
\paragraph{C.3au (Proved in the stated restricted scope): the rational ordinary top quotient of the complex Fermat quintic cone.}
Let
\[
 V=\{x_0^5+x_1^5+x_2^5+x_3^5=0\}\subset\mathbb P^3_{\mathbb C},
 \qquad L=\mathcal O_V(1),
\]
let $R$ be the strict henselization at the vertex of the local ring
of its affine cone, and put
$X=\operatorname{Spec}R\setminus\{\mathfrak m_R\}$.
Let $Z\to\operatorname{Spec}R$ be the regular cone resolution,
with exceptional divisor $i:V\hookrightarrow Z$ and projection
$p:Z\to V$. For ordinary algebraic equivalence on actual fixed
products, the asserted conclusion is
\begin{equation}\label{eq:C3au-answer}
 B^2(Z)_{\mathbb Q}\xrightarrow[\deg\circ i^!]{\ \sim\ }\mathbb Q,
 \qquad B^2(X)_{\mathbb Q}=0.
\end{equation}
This is a restricted statement about this specified complex surface
cone. In particular it is not an integral finiteness theorem, and
does not settle the rational top quotient for an arbitrary complex
quintic or an arbitrary isolated singularity.

\begin{proof}
We set $k_0=\overline{\mathbb Q}$ and use the displayed models
over $k_0$. Every differential called absolute below is relative
to $\mathbb Q$, equivalently to $k_0$ after base change.

\emph{1. The target is detected by its first infinitesimal neighbourhood.}
The full section ring is the displayed homogeneous hypersurface ring;
it is normal and generated in degree one. Thus the construction and
finite detection of C.3ab--ad apply. Write
\[
 V_a=\operatorname{Spec}_V
       (\mathcal O_V\oplus L\oplus\cdots\oplus L^{a-1}).
\]
The exact Adams-weight-two calculation in C.3ad, also recalled in
C.3ar, gives
\begin{equation}\label{eq:C3au-relative}
 K_0^{(2)}(V_a,V)_{\mathbb Q}
   =\bigoplus_{j=1}^{a-1}
       H^2(V,\Omega^1_{V/\mathbb Q}\otimes L^j).
\end{equation}
Adjunction gives $K_V=L$. Serre duality gives
\[
 H^2(V,L)=\mathbb C,
 \qquad H^2(V,L^j)=0\quad(j\ge2).
\]
The $k_0$-model splits the absolute differential sequence.
Kodaira--Akizuki--Nakano vanishing gives
$H^2(V,\Omega^1_{V/\mathbb C}\otimes L^j)=0$ for every $j>0$.
Consequently restriction to $V_2$ is an isomorphism on the
relative group in \eqref{eq:C3au-relative} for every $a\ge2$, and
that group is
\begin{equation}\label{eq:C3au-target}
       H^2(V,L)\otimes_{\mathbb C}\Omega^1_{\mathbb C/\mathbb Q}.
\end{equation}
All these equalities concern the underlying rational vector spaces.

There is a concrete normalization of this group useful below.
The integral Milnor Bloch cycle class and its symbol $d\log$
give a map
\[
 \mathrm{CH}^2(Z)\longrightarrow
                  H^2(Z,\Omega^2_{Z/\mathbb Q}).
\]
Restrict the differential sheaf by tensoring with $\mathcal O_V$.
The split first neighbourhood gives
\[
 \Omega^2_{Z/\mathbb Q}|_V
    =\Omega^2_{V/\mathbb Q}
          \oplus (L\otimes\Omega^1_{V/\mathbb Q}).
\]
Project to the conormal summand and then to its coefficient-field
component. Denote the resulting map by
\[
 d_{V,L}:\mathrm{CH}^2(Z)\longrightarrow
               H^2(V,L)\otimes\Omega^1_{\mathbb C/\mathbb Q}.
\]
The notation $|_V$ means tensor restriction, so the conormal
$dt$ term is retained. On the kernel of $i^!$, this is the
relative identification \eqref{eq:C3au-target} with coefficient
one: locally the horizontal cyclic representative $t\omega$
has differential $dt\wedge\omega+t\,d\omega$, and its tensor
restriction is $dt\wedge\omega$. This is precisely the symbol
normalization checked in C.3ar, Step~3, using
\citet[Lemma 2.2, Proposition 2.3 and the proof of Proposition 6.1]{D_BlochEsnaultKerz2014}.
The fixed-level group calculation itself remains the Loday and
Corti\~nas--Haesemeyer--Weibel calculation of C.3ad.

In particular, finite detection on $K_0(Z)$ and the rational
Chern character on the regular finite pointed-\'etale models give
\begin{equation}\label{eq:C3au-detect}
 i^!z=0\ \hbox{and}\ d_{V,L}(z)=0
       \quad\Longrightarrow\quad z=0
           \quad\bigl(z\in\mathrm{CH}^2(Z)_{\mathbb Q}\bigr).
\end{equation}
No assertion about algebraization of a general formal $K_0$
class is used in this implication.

\emph{2. A specified degree-five cover by a blown-up product of curves.}
Consider the two smooth plane quintics over $k_0$
\[
 C=\{a^5+b^5+c^5=0\}\subset\mathbb P^2,
 \qquad
 B=\{v^5=t^5+w^5\}\subset\mathbb P^2.
\]
Both have genus six and canonical bundle $\mathcal O(2)$.
On $C\times B$ the four sections
\[
                  (tc,\ wc,\ va,\ vb)
\]
of $M=\mathcal O_C(1)\boxtimes\mathcal O_B(1)$ satisfy
the equation of $V$. Their base scheme consists of the twenty-five
distinct points $c=v=0$. At each such point, $c,v$ are regular
parameters, the remaining relevant coordinates are units, and the
base ideal is exactly $(c,v)$. A single blow-up at this reduced
set resolves the map. Write
\[
 \pi:T=\operatorname{Bl}_{\{c=v=0\}}(C\times B)\longrightarrow C\times B,
 \qquad f:T\longrightarrow V,
 \qquad E=\sum_{\nu=1}^{25}E_\nu.
\]
Thus $T$ is smooth and projective, $f$ is proper, and
\begin{equation}\label{eq:C3au-line}
             L':=f^*L=\pi^*M\otimes\mathcal O_T(-E).
\end{equation}
On the open set with $w=c=1$, a general image point determines
$t=x_0/x_1$, and $v$ has exactly five possible nonzero values
with $v^5=t^5+1$. Each determines
\[
                    a=\frac{x_2}{v x_1},\qquad
                    b=\frac{x_3}{v x_1}.
\]
These are exactly the five inverse points. Hence $f$ is dominant
and generically finite of degree five, and is generically \'etale.
We do not call $L'$ ample; none of the following steps requires it.

Let
\[
              Z'=Z\times_V T,\qquad F:Z'\longrightarrow Z.
\]
This is the strict-local base change of the total space of
$(L')^{-1}$, so it is regular and proper over $R$.
Its exceptional divisor is $i':T\hookrightarrow Z'$, and its
formal neighbourhoods are the split truncated algebras of powers
of $L'$. The equality $\mathfrak m_R\mathcal O_Z=\mathcal O_Z(-V)$
pulls back to
$\mathfrak m_R\mathcal O_{Z'}=\mathcal O_{Z'}(-T)$.
In particular, ``modulo $\mathfrak m_R^2$'' is exactly the first
infinitesimal neighbourhood of $T$ used below.

\emph{3. Any chosen Picard first jet on $Z'$ is realized by an actual line bundle.}
Let $a\in H^1(T,L')$. On the split formal completion of $Z'$
define a line bundle by the transition functions
$\exp(t a_{\alpha\beta})$, where $a_{\alpha\beta}$ represents
$a$. The formal logarithm has only homogeneous degree one.
The cocycle identity holds in every order. Formal existence for
the proper scheme over $\widehat R$ gives an actual line bundle
$\widehat A$ on $Z'_{\widehat R}$
\citep[Theorem 30.27.1, Tag 088E]{rootStacksExistence}.
Invertibility follows from the invertible special fibre by
Nakayama and properness: the closed locus where the algebraized
coherent sheaf is not a line bundle cannot have nonempty proper
image in a local base while avoiding its closed point.

We only descend the specified finite jet. Excellence makes
$R\to\widehat R$ regular. Popescu's theorem writes $\widehat R$
as a filtered colimit of smooth finitely presented $R$-algebras
$A_\nu$ \citep[Theorem 16.12.1, Tag 07GC]{cyStacksPicardJet}.
The line bundle $\widehat A$ descends to a line bundle on
$Z'_{A_\nu}$ for some $\nu$. Indeed a finite affine cover, its
finite transition functions, their inverses and cocycle relations
are finite presentation data; the standard limit theorem for
perfect complexes gives the same conclusion after enlarging
$\nu$ \citep[Lemma 36.29.3, Tag 09RF]{cyStacksPicardJet}.
The map $A_\nu\to\widehat R$ induces $A_\nu\to R/\mathfrak m_R^2$.
Since $(R,\mathfrak m_R^2)$ is a henselian pair, smooth lifting
gives an $R$-algebra map $A_\nu\to R$ inducing that same map
modulo $\mathfrak m_R^2$
\citep[Lemma 15.13.3, Tag 0H74]{cyStacksPicardJet}.
Pulling back gives a line bundle $A_a$ on the actual $Z'$ with
\begin{equation}\label{eq:C3au-jet}
          A_a|_T=\mathcal O_T,
          \qquad \log(A_a|_{T_2})=t a.
\end{equation}
We make no assertion that its higher formal coefficients agree
with those of $\widehat A$. They are irrelevant by
\eqref{eq:C3au-detect} on the target. In particular, this argument
does not use Serre vanishing for $L'$ and does not algebraize
an arbitrary inverse system of $K_0$ classes.

\emph{4. Actual products span every coefficient on the source.}
Let $q_B:T\to B$ be the projection through $\pi$. For
$\lambda=\mathcal O_B(q-o)$, where $o\in B(k_0)$ and
$q\in B(\mathbb C)$, define
\[
             \beta(a,\lambda)
               =c_1(A_a)c_1(q_B^*\lambda)
                         \in\mathrm{CH}^2(Z').
\]
Here and below $q_B$ also denotes its composite with $Z'\to T$.
This is ordinary algebraically trivial: keep $A_a$ fixed and
use the actual diagonal-minus-$o$ divisor on $B\times B$ as
the parameter family. Its restriction to $T$ is zero.
As in C.3ar, the pure coefficient-field component of its
conormal $d\log$ class is
\begin{equation}\label{eq:C3au-product}
       a\smile q_B^*\epsilon(\lambda)
        \in H^2(T,L')\otimes\Omega^1_{\mathbb C/\mathbb Q},
\end{equation}
with coefficient one, up to one fixed harmless cup-product sign.
The relative first Chern class of the degree-zero $\lambda$ is
zero, so this formula omits no relative differential contribution.

We check its span explicitly. The elementary blow-up calculation
$R\pi_*\mathcal O_T(-E)=\mathcal I_{\{c=v=0\}}$ has no higher
direct images. This follows either on the two standard affine
charts at each point or from the exceptional $\mathbb P^1$ and
its nonnegative conormal powers. The sequence of the reduced
finite base scheme and the projection formula give
\begin{equation}\label{eq:C3au-coherent}
 H^1(T,L')\twoheadrightarrow H^1(C\times B,M),
 \qquad H^2(T,L')\simeq H^2(C\times B,M).
\end{equation}
These maps commute with cup product by classes pulled back from
$C\times B$. In particular
\[
 H^2(T,L')\simeq
 H^1(C,\mathcal O_C(1))\otimes H^1(B,\mathcal O_B(1)).
\]
Choose a nonzero $\sigma\in H^0(B,\mathcal O_B(1))$.
Cup product by $\sigma$ is surjective:
\[
       H^1(B,\mathcal O_B)\xrightarrow{\,\sigma\,}
                          H^1(B,\mathcal O_B(1)).
\]
Indeed its Serre-dual map
\[
       H^0(B,\mathcal O_B(1))\longrightarrow H^0(B,\mathcal O_B(2))
\]
is injective.

For complete control of coefficient forms use
$u=(t/w)^3:B\to\mathbb P^1$. This has degree fifteen, greater
than $2g(B)-2=10$. For any transcendental $b\in\mathbb C$ its
fibre has fifteen distinct unramified points. The residue formula
of C.3ao and C.3ar gives
\[
 \left\langle\epsilon\bigl(\mathcal O_B(q-o)\bigr),\omega\right\rangle
       =-q^*\omega
       =-\left(\frac{\omega}{du}\right)(q)\,db.
\]
Evaluation at these points spans the dual of $H^0(B,K_B)$:
a nonzero canonical form cannot vanish at fifteen points.
Hence the coefficient vectors of these actual Picard classes
span $H^1(B,\mathcal O_B)\otimes db$.
Lift any $\eta\otimes\sigma$, with
$\eta\in H^1(C,\mathcal O_C(1))$, through
\eqref{eq:C3au-coherent} to an $a\in H^1(T,L')$.
Formula \eqref{eq:C3au-product} and these fibre evaluations
therefore realize every element of
\begin{equation}\label{eq:C3au-full-span}
                    H^2(T,L')\otimes\Omega^1_{\mathbb C/\mathbb Q}
\end{equation}
by a finite sum of actual ordinary products $\beta(a,\lambda)$.
As in C.3ar, every complex scalar appearing in this linear
algebra is absorbed into the chosen $a$ before constructing
$A_a$. We do not multiply a Chow class by an arbitrary complex
number.

\emph{5. The exact first-jet pushforward.}
Let
\[
 \operatorname{tr}_{f,L}:H^2(T,L')\longrightarrow H^2(V,L)
\]
be the Serre-dual map to
\[
 f^*:H^0(V,K_V\otimes L^{-1})
                      \longrightarrow H^0(T,K_T\otimes(L')^{-1}).
\]
This is the ordinary differential pullback followed by
$f^*L=L'$. It is nonzero, since $f$ is generically \'etale.
Its source is one-dimensional. Accordingly
$\operatorname{tr}_{f,L}$ is surjective.

For the cycles in Step~4 the compatibility we need is
\begin{equation}\label{eq:C3au-trace}
 d_{V,L}\bigl(F_*\beta(a,\lambda)\bigr)
   =(\operatorname{tr}_{f,L}\otimes1)
           \bigl(a\smile q_B^*\epsilon(\lambda)\bigr).
\end{equation}
We give the normalization and finite-stage justification, since
compatibility with an arbitrary formal pushforward would not be
a substitute for this assertion.

At a smooth finite pointed-\'etale model the Hodge cycle class
is the fundamental class defined by its local residue symbol
\[
                   \left[\frac{dg_1\wedge dg_2}{g_1,g_2}\right]
\]
at a regular generic point of a codimension-two cycle.
Its extension across singularities is defined by normalization
and residue trace. This agrees with the Milnor $d\log$ class:
both send the symbol $\{g_1,g_2\}$ to the displayed class,
and the Gersten-residue comparison identifies their differentials
\citep[Section~4, pp.~559--560]{cyBraunling2013}.
The residue trace sends a proper dominant map of cycle supports
of equal dimension to its generic degree times the target
fundamental class; a lower-dimensional image contributes zero
in this degree. These are precisely the coefficients in proper
Chow pushforward. The exact primary statements are
\citet[Section~2.4, pp.~31--32, and Section~3.1, Proposition~2,
p.~41]{cyElZein1978}, with the finite trace normalization also
written in \citet[Proposition~2.1.2, pp.~15--16]{cyElZein1974}.

For the present products the restriction calculation can be made
with supports meeting $T$ properly. Represent $q_B^*\lambda$
by the difference of the two pulled-back fibre divisors. Choose
a rational section of $A_a$ which is a unit both at the generic
points of their intersections with $T$ and at the generic
points of the fibre divisors themselves. Such a choice exists by
trivializing the line bundle on the semilocal ring of this
finite set of points on a finite quasi-projective model.
Its divisor intersects the two fibre divisors in a cycle
representing $\beta(a,\lambda)$ with no component in $T$.
The intersection with $T$ is therefore a finite set.
Normalizing those curve supports reduces the calculation at
that set to traces of finite maps of curve branches; intersection
multiplicities are retained by the residue formula. Thus the
argument does not require a pullback of a fraction whose whole
support lies in the exceptional divisor.

We spell out why this trace also commutes with the tensor
conormal projection. First choose a $\mathbb Q$-derivation
$\theta:\mathbb C\to\mathbb C$. The $k_0$-models and their
\'etale localizations extend $\theta$ while fixing the geometric
coordinates. Apply $-\iota_\theta$ to the absolute degree-two
Hodge class and project its remaining differential to
$\Omega^1_{Z/\mathbb C}$. Denote the result by $\xi_\theta$.
Its conormal coefficient is precisely the contraction of
$d_{V,L}(z)$ by $\theta$: the sign is fixed by
$\iota_\theta(dt\wedge\omega_{\mathrm{coeff}})
       =-dt\,\theta(\omega_{\mathrm{coeff}})$.

For $\eta\in H^0(V,K_V\otimes L^{-1})$, write $\omega$ for
its top-form representative in a local frame of $L$, and
write $t$ for the corresponding normal coordinate. The
meromorphic relative two-form $\omega/t$ is independent of
the frame and has a simple pole along $V$. Wedge with this
form and take the simple-pole residue. Locally
\[
 \operatorname{Res}_V
       \left((dt\,b+\xi_{\mathrm{horizontal}})
                           \wedge\frac{\omega}{t}\right)
                 =b\omega.
\]
The horizontal wedge vanishes because it is pulled back from
three relative differentials on the surface $V$.
Consequently Serre trace of this residue is exactly the
pairing $\langle (1\otimes\theta)d_{V,L}(z),\eta\rangle$.
This formula involves only relative forms after contraction;
there is no implicit scalar-valued pairing of an uncontracted
absolute field differential.

Pullback by $F$ sends $\omega/t$ to $f^*\omega/t$. Proper
Hodge trace is adjoint to this pullback, including for the
supported residue pairing. The latter assertion can either
be read directly in the supported duality construction of
El Zein, Section~2.4, or checked on the residue representatives
above: multiply a local fundamental-class symbol by
$\omega/t$, apply finite residue trace on its normalized curve
branches, and sum. The finite trace formula commutes with
this multiplication and with residue; its branch degrees are
exactly those of proper cycle pushforward. This is the
residue construction on pp.~32--33 of
\citet{cyElZein1978}. It proves the required compatibility with
the conormal coefficient, including ramification and singular
supports. No ramification locus is discarded.

Since $f$ and $L,L'$ are defined over $k_0$, contraction by
$\theta$ commutes with this trace. The resulting source
pairing is
\[
 \left\langle
  (1\otimes\theta)\bigl(a\smile q_B^*\epsilon(\lambda)\bigr),
                  f^*\eta\right\rangle.
\]
By definition this is the pairing with
$\operatorname{tr}_{f,L}$ in \eqref{eq:C3au-trace}.
Derivations separate the coefficient-field differentials,
and $\eta$ spans the Serre dual of $H^2(V,L)$, so these
pairing identities prove that equation with its stated
normalization.

All cycles, rational equivalences and line bundles in this
calculation involve finitely many coefficients and descend to
one smooth finite-type coefficient model and one finite
pointed-\'etale neighbourhood. The local residue identities
are compatible with \'etale localization and scalar extension,
so their filtered limit proves \eqref{eq:C3au-trace} for the
actual strict-local schemes. Equivalently one may evaluate the
coefficient-field component against any derivation of its
finitely generated coefficient field, perform this same
residue calculation there, and pass to the filtered union.
This uses finite data and actual proper pushforward throughout.

The Cartier square is Tor-independent and $F^{-1}(V)=T$
with multiplicity one. Refined pullback therefore also gives
\[
                   i^!F_*\beta=f_*i'^!\beta=0.
\]
Proper pushforward preserves the ordinary family from Step~4:
apply the actual map $F\times1_B$ to that family.

\emph{6. Elimination and localization.}
Let $z\in\mathrm{CH}^2(Z)_{\mathbb Q}$ satisfy $i^!z=0$.
By \eqref{eq:C3au-full-span}, surjectivity of
$\operatorname{tr}_{f,L}$ and \eqref{eq:C3au-trace}, there is
a finite sum of actual ordinary cycles $\beta_\nu$ on $Z'$
such that
\[
          d_{V,L}\left(z-\sum_\nu F_*\beta_\nu\right)=0.
\]
The expression also has zero $i^!$, so
\eqref{eq:C3au-detect} makes it zero rationally. Thus every
element of $\ker i^!$ is ordinarily algebraically trivial.
For an arbitrary class subtract $p^*i^!z$ first.
The ordinary zero-cycle quotient of the smooth connected
projective $V$ is degree, giving the first isomorphism in
\eqref{eq:C3au-answer}. Finally localization and
$i^!i_*=-c_1(L)$ show that the exceptional image contains
a nonzero multiple of this degree generator: take any
integral curve of positive $L$-degree on $V$.
The rational ordinary quotient on $X$ is therefore zero.
\end{proof}
